\chapter{The AMOC: Overturning, Hosing, and Tipping}
\label{chap:ch17_amoc_tipping}

The Atlantic Meridional Overturning Circulation (AMOC) carries warm, salty water poleward in
the upper Atlantic and returns cold, dense water at depth. It is a textbook example of a climate
\emph{tipping element}: theory and a hierarchy of models suggest the overturning can collapse
abruptly and \emph{irreversibly} once a freshwater perturbation exceeds a critical threshold,
because of the salt-advection feedback first identified by \citet{stommel1961}. This chapter
uses \model{} as a laboratory to reproduce that bifurcation, and --- exploiting
differentiability --- to \emph{calibrate} where it occurs. The governing equations of the
overturning closure are given in Chapter~\ref{chap:ch08_ocean_overturning}; here we focus on the
experiments and results.

\section{The salt-advection feedback}

Write the overturning strength $q$ as an increasing function of the meridional density contrast
between the subpolar and subtropical Atlantic. A stronger $q$ imports more salt to the subpolar
basin, raising its density, which strengthens $q$ further: a positive feedback. A subpolar
freshwater input (``hosing'') opposes it. When the feedback loop gain exceeds unity the system
becomes \emph{bistable} --- an ``on'' (vigorous) and an ``off'' (collapsed) state coexist over a
finite range of forcing --- and the transition between them is a saddle-node bifurcation with
hysteresis. In \model{} the loop gain is set by the tunable parameter \code{haline\_gain}
(Chapter~\ref{chap:ch08_ocean_overturning}).

\section{The hosing experiment}

We add a freshwater flux $F$ (in Sverdrups, $1\Sv=10^6~\mathrm{m^3\,s^{-1}}$) over the
$45$--$65^\circ$N Atlantic surface and record the AMOC strength at $26.5^\circ$N. Two protocols
are used:

\begin{description}[leftmargin=0pt,style=nextline]
  \item[Quasi-static sweep] Ramp $F$ slowly up ($0\!\to\!F_{\max}$) then down
    ($F_{\max}\!\to\!0$), holding each level to near-equilibrium. An open loop (the down-leg
    below the up-leg) signals hysteresis. \emph{Caveat:} near the fold the collapse timescale is
    multi-decadal, so finite holds give a \emph{rate-dependent} loop that can be mistaken for, or
    hide, true bistability.
    \codref{experiments/run\_amoc\_hosing.py}{quasi-static hosing sweep}
  \item[Fixed-$F$ bistability test] The rigorous protocol: at a single fixed $F$, integrate one
    run started from an \emph{on}-state and one from an \emph{off}-state for $\sim100$~years. If
    the two converge, the system is monostable at that $F$; if they remain separated, it is
    genuinely bistable. This is immune to rate-dependence.
    \codref{experiments/run\_amoc\_bistability.py}{fixed-$F$ on/off initial-condition test}
\end{description}

\section{A verified bistable window}

With the calibrated configuration \code{haline\_gain}$=6$, \code{thc\_k\_vel}$=6\times10^{-5}$,
and a convection-layer contrast depth of $300$~m, the fixed-$F$ test gives the result in
Table~\ref{tab:amoc_bist} and Figure~\ref{fig:amoc_bist}. At $F=0.46,\,0.57,\,0.69\Sv$ the
on-state and off-state branches remain $8$--$9\Sv$ apart after $100$~years --- a separation that
is statistically significant (autocorrelation-corrected $z=2.6$--$4.4$) and persistent (the
off-branch does not recover). Below $F\approx0.3\Sv$ an off-state initial condition recovers to
the upper branch; above $F\approx0.8\Sv$ an on-state initial condition collapses. The
\emph{hysteresis window} is therefore approximately
\[
  F_{\mathrm{lower}} \approx 0.38\Sv \;<\; F \;<\; F_{\mathrm{upper}} \approx 0.75\Sv,
\]
centered near $0.6\Sv$ --- consistent with the freshwater-transport scale of a $\sim15$--$20\Sv$
AMOC and with the collapse thresholds reported across the model hierarchy
\citep{rahmstorf2005}.

\begin{table}[h]
\centering
\caption{Equilibrated AMOC (last-40-yr mean) at $26.5^\circ$N from on-state vs.\ off-state
initial conditions, held at fixed hosing $F$ for $100$~years.}
\label{tab:amoc_bist}
\begin{tabular}{cccc}
\toprule
$F$ (Sv) & ON-branch (Sv) & OFF-branch (Sv) & separation $z$ \\
\midrule
$\le 0.30$ & \multicolumn{2}{c}{off-IC recovers to ON} & --- \\
$0.46$ & $\approx 20$ & $\approx 12$ & $2.6$ \\
$0.57$ & $\approx 18$ & $\approx 10$ & $3.1$ \\
$0.69$ & $\approx 15$ & $\approx 6$  & $4.4$ \\
$\ge 0.80$ & \multicolumn{2}{c}{on-IC collapses to OFF} & --- \\
\bottomrule
\end{tabular}
\end{table}

\begin{figure}[h]
\centering
\includegraphics[width=\linewidth]{amoc_bistability}
\caption{The AMOC hysteresis window. (a) On-state (blue) and off-state (red) equilibrated AMOC
versus hosing $F$; the branches stay separated across the window. (b) At $F=0.57\Sv$ the two
initial conditions settle on different branches and do not reconverge over $100$~years ---
the signature of genuine bistability.}
\label{fig:amoc_bist}
\end{figure}

\paragraph{An honest caveat.} The branches are \emph{noisy}: the AMOC oscillates by
$\sim\!\pm10\Sv$ (a relaxation oscillation intrinsic to the strong feedback in this interim
box-style closure). The bifurcation is genuine --- established by the statistically-significant,
non-recovering initial-condition dependence --- but it is not a clean, low-variance loop. A clean
dynamical bifurcation requires the prognostic-momentum ocean core of
Chapter~\ref{chap:ch07_ocean_prognostic_core}; the closure used here \emph{works around}, rather
than solves, the fact that the diagnostic-velocity ocean has $\partial(\mathrm{AMOC})/\partial\rho
\approx 0$.

\section{Calibrating the tipping point by differentiation}

A tipping threshold is a bifurcation, and therein lies a subtlety for differentiable modeling:
one \emph{cannot} obtain $\partial F_{\mathrm{crit}}/\partial\boldsymbol{\theta}$ by
backpropagating through the hosing sweep, because the trajectory's sensitivity \emph{diverges} at
the fold (that divergence \emph{is} the bifurcation). The correct tool is \emph{fold
continuation}: reduce the closure to its governing Stommel box, where the steady-state condition
$F = R(s;\boldsymbol{\theta})$ has a fold at $\partial R/\partial s=0$, so that
$F_{\mathrm{crit}}(\boldsymbol{\theta}) = \max_s R(s;\boldsymbol{\theta})$. By the envelope
theorem the gradient through this maximum is exact, and a few Newton steps invert
$F_{\mathrm{crit}}(\boldsymbol{\theta})=F^\star$ for the parameters in milliseconds --- the
differentiable replacement for a brute-force parameter grid.
\codref{experiments/calibrate\_amoc\_fold.py}{AD fold-continuation calibration}

\noindent
The gradient ranking this yields,
$\partial F_{\mathrm{crit}}/\partial(\code{k\_vel}) \gg
 \partial F_{\mathrm{crit}}/\partial(\code{haline\_gain})$,
shows that the AMOC on-state \emph{strength} is the dominant control on the tipping threshold ---
a non-obvious result that the grid search would have taken hours to reveal.

\section{Reproducing this chapter}

\begin{lstlisting}[language=bash,caption={Hosing sweep, bistability test, and figure.}]
# quasi-static hysteresis sweep (self-contained from WOA; auto-resumes):
python experiments/run_amoc_hosing.py --spinup-years 20 --fmax 0.8 --nsteps 8 \
    --hold-years 15 --haline-gain 6 --k-vel 6e-5 --contrast-depth 300

# rigorous fixed-F bistability test (run once per branch IC, then compare means):
python experiments/run_amoc_bistability.py --ckpt <on_state>  --hosing-sv 0.57 \
    --haline-gain 6 --k-vel 6e-5 --contrast-depth 300 --years 100 --outdir out/on
python experiments/run_amoc_bistability.py --ckpt <off_state> --hosing-sv 0.57 \
    --haline-gain 6 --k-vel 6e-5 --contrast-depth 300 --years 100 --outdir out/off

# differentiable calibration of the threshold, and the figure:
python experiments/calibrate_amoc_fold.py --target 0.6
python experiments/plot_amoc_bistability.py --root out --out docs/figures/amoc_bistability.pdf
\end{lstlisting}

\begin{exercise}
Re-run the fixed-$F$ test at $F=0.57\Sv$ but lower \code{haline\_gain} to $4$. Is the system still
bistable? Relate your answer to the loop-gain condition for bistability.
\end{exercise}

\begin{exercise}
Using \code{calibrate\_amoc\_fold.py}, find the \code{(haline\_gain, k\_vel)} pair that places the
collapse threshold at exactly $F_{\mathrm{crit}}=0.5\Sv$. Verify with one fixed-$F$ run near your
predicted threshold.
\end{exercise}

\begin{exercise}[Differentiability]
Explain in one paragraph why $\partial F_{\mathrm{crit}}/\partial\boldsymbol{\theta}$ cannot be
computed by automatic differentiation \emph{through} a hosing sweep that crosses the fold, and how
fold continuation circumvents the problem.
\end{exercise}
